\section{Related work}
\label{sec:related}

\paragraph{Tactile representations and evaluation.}
Self-supervised tactile learning aims to extract transferable information about contact from large image collections, as in Sparsh~\cite{higuera2024sparsh}, T3~\cite{t3fota2024}, UniTouch~\cite{unitouch2024}, and AnyTouch~\cite{anytouch2025,anytouch2026two}. Other work learns representations across modalities or views of contact~\cite{tvl2024,tactspace2026}. Evaluating these representations requires separating contact information from cues associated with data acquisition. RCT demonstrates sequence leakage under frame-level splits of TVL~\cite{rct2026}, while a recent material-recognition benchmark on HCT shares recording sessions between training and test~\cite{mohsan2026language}. These observations motivate evaluation beyond randomly held-out images. Hardware introduces a further concern: even an object-level split can preserve the same sensor appearance across training and test. We examine whether classifiers exploit that appearance when material classes are associated with particular hardware configurations, and evaluate recognition after changes of unit or gel over four months.

\paragraph{Hardware variation in tactile sensing.}
Existing tactile pipelines already account for sensor appearance in several ways. Background subtraction removes an unloaded-gel reference image~\cite{higuera2024sparsh,sitr2025,anytouch2026two,htt2026}; GelSlim periodically refreshes this reference~\cite{gelslim2018}; and sensor-type conditioning provides information about the acquisition device~\cite{unitouch2024}. Studies of gel durability also characterize changes in unloaded images with use~\cite{davis2025resilience,cottrell2026durable}. Transfer methods address differences more directly through calibration contacts and simulation~\cite{sitr2025}, labelled target contacts~\cite{bideta2026}, target-background-conditioned translation~\cite{genforce2026,castano2026forcemap,split2026}, or alignment across sensor types~\cite{cttp2024,htt2026}. FEATS and FeelAnyForce document degradation in force estimation across units or gels~\cite{feats2025,feelanyforce2024}, and TacVerse evaluates multiple sensors with chronological splits~\cite{tacverse2026}. This work establishes the importance of hardware variation. Our focus is the persistent image structure associated with a particular unit and gel, how classifiers use it, and how its contribution can be estimated from ordinary unlabelled contacts without additional calibration presses or unloaded-gel frames.

\paragraph{Acquisition shortcuts.}
Differences between acquisition systems become classification shortcuts when they correlate with labels~\cite{geirhos2020shortcut}. Examples include hospital-specific cues in chest radiographs~\cite{zech2018hospital}, acquisition differences in COVID datasets~\cite{degrave2021covid}, and dataset-specific appearance in natural images~\cite{torralba2011bias,liuhe2025}. Tactile datasets can create an analogous association when materials are recorded in separate sessions on different hardware configurations. Our contribution is to measure the hardware signature over time and test its effect on recognition, including controlled class--hardware associations in public data. This distinguishes a classifier's reliance on hardware from a general decline under distribution shift.

\paragraph{Estimating acquisition effects from images.}
Estimating a shared acquisition component from multiple images is well established. Microscopy estimates illumination fields by averaging images or fitting low-rank structure~\cite{singh2014illumination,peng2017basic}; infrared imaging estimates fixed-pattern noise~\cite{harris1999nuc}; and camera identification recovers sensor-specific signatures~\cite{lukas2006prnu}. Multi-site MRI similarly models scanner-dependent measurement effects~\cite{fortin2017combat}. We apply this principle to a robot's contact stream: observations of different surfaces provide a running estimate of the sensor's hardware state. The correction removes only its smooth component, since fine-scale averages can retain object texture as well as hardware structure. The estimate is restarted after detected hardware changes. The contribution is therefore a state-aware tactile reading grounded in measured hardware persistence, rather than a new general-purpose flat-field estimator.

\paragraph{Adaptation from a robot's contact stream.}
Unsupervised domain adaptation uses labelled source data and unlabelled target observations to address distribution shift through feature alignment~\cite{ganin2016dann,long2018cdan}, class-confusion reduction~\cite{jin2020mcc}, or regularization~\cite{rangwani2022sdat,hoyer2023mic}. Test-time adaptation instead updates statistics or model parameters during inference~\cite{li2018adabn,wang2021tent,niu2023sar,lee2024deyo}, with dedicated methods addressing temporally correlated inputs~\cite{gong2022note,boudiaf2022lame,yuan2023rotta}. This distinction matters for tactile sensing: consecutive contacts may involve one material, so batch statistics can absorb class information rather than isolate hardware variation~\cite{burns2021bn}. Our correction explicitly estimates a persistent image component without updating the classifier. It is complementary to feature adaptation, for which we use the same contact stream as the target domain. Unlike FDA, which transfers low-frequency appearance into source training images~\cite{fda2020}, our correction adjusts test images relative to the estimated hardware state. We evaluate both correction alone and its addition to existing adaptation methods.